\documentclass[a4paper,12pt]{article}

\usepackage[T1]{fontenc}
\usepackage[italian]{babel}
\usepackage{graphicx}
\usepackage{geometry}
\usepackage{tabularx}
\usepackage[table]{xcolor}
\usepackage[most]{tcolorbox}
\usepackage{fancyhdr}
\usepackage[hidelinks]{hyperref}

\newcommand{\groupmail}{alittlebyte19@gmail.com}
\newcommand{\grouplogo}{../../.github/site-src/assets/logo_alittlebyte.png}
\newcommand{\groupsymbol}{../../.github/site-src/assets/solo_logo.png}

\geometry{
    left=2.5cm,
    right=2.5cm,
    top=2.5cm,
    bottom=2.5cm,
    headheight=331pt,
    includeheadfoot
}

\pagestyle{fancy}
\fancyhf{}

\renewcommand{\headrulewidth}{0pt}
\fancyhead{}

\renewcommand{\footrulewidth}{0.4pt}
\fancyfoot[L]{\includegraphics[height=0.6cm]{\groupsymbol}\hspace{0.45em}aLittleByte}
\fancyfoot[R]{Analisi dei Capitolati}

\fancypagestyle{firstpage}{
    \fancyhf{}
    \renewcommand{\headrulewidth}{0pt}
    \renewcommand{\footrulewidth}{0.4pt}
    \fancyhead[C]{%
        \makebox[\textwidth][c]{%
            \begin{tabular}{c}
                \includegraphics[height=10cm]{\grouplogo}\\[1ex]
                \small\texttt{\groupmail}
            \end{tabular}%
        }
    }
    \fancyfoot[L]{\includegraphics[height=0.6cm]{\groupsymbol}\hspace{0.45em}aLittleByte}
    \fancyfoot[R]{Analisi dei Capitolati}
}

\begin{document}

\thispagestyle{firstpage}

\vspace*{0.5cm}

\begin{center}
    {\LARGE \textbf{Analisi dei Capitolati}}
\end{center}

\vspace{1.5cm}

\begin{center}
    \renewcommand{\arraystretch}{1.3}
    \begin{tcolorbox}[
        width=0.92\textwidth,
        colback=gray!6,
        colframe=gray!45,
        boxrule=0.5pt,
        arc=2mm,
        boxsep=0pt,
        left=8pt, right=8pt, top=8pt, bottom=8pt
    ]
        \begin{tabularx}{\linewidth}{>{\bfseries}p{3.5cm} X}
            Gruppo      & aLittleByte (gruppo 19) \\
            Redattori   & Angelica Gastaldello, Diego Belluz \\
            Verificatori& Leonardo Soligo, Elenora Bellet\\
            Destinatari & Prof. Tullio Vardanega, Prof. Riccardo Cardin \\
            Data        & 20/03/2026 \\
        \end{tabularx}
    \end{tcolorbox}
\end{center}

\newpage

\newgeometry{left=2.5cm,right=2.5cm,top=2.5cm,bottom=2.5cm,includefoot}

\begin{center}
    {\Large \textbf{Registro delle Modifiche}}
\end{center}

\vspace{0.35cm}

\renewcommand{\arraystretch}{1.35}
\rowcolors{2}{gray!8}{white}
\begin{center}
    {\footnotesize
    \begin{tabularx}{\textwidth}{|p{1.4cm}|p{1.8cm}|X|p{2.4cm}|p{2.3cm}|p{2.3cm}|}
        \hline
        \rowcolor{gray!20}
        \textbf{Versione} & \textbf{Data} & \textbf{Descrizione} & \textbf{Redattore} & \textbf{Verificatore} & \textbf{Approvatore} \\
        \hline
        1.0.0 & 20/03/2026 & Approvazione documento &  &  & Eleonora Bellet, Leonardo Soligo \\
        \hline
        0.5.2 & 19/03/2026 & Aggiunti link verbali & Angelica Gastaldello & Diego Belluz & - \\
        \hline
        0.5.1 & 18/03/2026 & Aggiornamento post-incontri con i proponenti & Angelica Gastaldello & Leonardo Soligo & - \\
        \hline
        0.5.0 & 12/03/2026 & Stesura analisi C5 & Angelica Gastaldello & Leonardo Soligo & - \\
        \hline
        0.4.0 & 11/03/2026 & Stesura analisi capitolato C6 & Diego Belluz & Elenora Bellet & - \\
        \hline
        0.3.0 & 10/03/2026 & Stesura analisi capitolato C4 & Diego Belluz & Elenora Bellet & - \\
        \hline
        0.2.0 & 10/03/2026 & Stesura analisi C2 & Angelica Gastaldello & Leonardo Soligo & - \\
        \hline
        0.1.0 & 09/03/2026 & Stesura analisi C1 & Angelica Gastaldello & Leonardo Soligo & - \\
        \hline
        0.0.2 & 08/03/2026 & Introduzione e Criteri & Angelica Gastaldello & Elenora Bellet & - \\
        \hline
        0.0.1 & 07/03/2026 & Struttura del documento & Angelica Gastaldello & Elenora Bellet & - \\
        \hline
    \end{tabularx}
    }
\end{center}
\newpage

\begin{center} 
    {\Large \textbf{Indice}}
\end{center}

\vspace{0.6cm}

\tableofcontents

\newpage

\pagenumbering{arabic}
\setcounter{page}{4}
\fancyfoot[R]{Analisi dei Capitolati\hspace{0.8em}\thepage}

\section{Introduzione}
\subsection{Finalità del documento}

Il presente documento ha lo scopo di illustrare l'analisi e la valutazione dei capitolati d'appalto proposti per l'anno accademico 2025/2026. Al suo interno viene fornita una disamina dettagliata del processo decisionale che ha portato il gruppo aLittleByte alla scelta finale.

Nello specifico, il documento approfondisce le motivazioni tecniche e organizzative alla base della selezione del capitolato C5 - Nexum, proposto dall'azienda Eggon S.r.l. Il progetto prevede lo sviluppo di moduli innovativi di AI Generativa per una piattaforma HR evoluta, integrando tecnologie cloud all'avanguardia. Per garantire la completezza della trattazione, viene inoltre riportata una valutazione comparativa degli altri capitolati disponibili, evidenziando i criteri di confronto utilizzati dal team.

\section{Criteri e metodologia}

La scelta del capitolato da parte del gruppo ha seguito un approccio analitico e strutturato, basato sulla definizione di precisi parametri di valutazione e sull'applicazione di un processo decisionale articolato in più fasi.

\subsection{Fattori considerati nella scelta}

Per valutare e confrontare oggettivamente le diverse proposte aziendali, il team ha individuato i seguenti criteri fondamentali:

\begin{itemize}
    \item \textbf{Complessità e fattibilità:} si è cercato un attento bilanciamento tra il livello di sfida tecnica (in termini di architettura e integrazione di sistemi) e la reale sostenibilità del carico di lavoro. L'obiettivo era garantire un prodotto finale di qualità nel pieno rispetto del monte ore a disposizione.
    \item \textbf{Interesse:} è stato dato un peso rilevante alla curiosità, alla motivazione e al coinvolgimento del gruppo verso il dominio applicativo del progetto, considerando la motivazione condivisa un elemento determinante.
    \item \textbf{Chiarezza delle tecnologie:} si è valutata la trasparenza e la precisione dei requisiti tecnologici richiesti, un aspetto cruciale per stimare preventivamente la curva di apprendimento necessaria e agevolare la futura divisione dei compiti.
    \item \textbf{Disponibilità del proponente:} è stata considerata essenziale la presenza di referenti tecnici aziendali pronti a fornire un supporto continuo, al fine di mitigare i rischi e chiarire tempestivamente eventuali dubbi in fase di sviluppo.
\end{itemize}

\subsection{Processo decisionale}

L'applicazione dei criteri appena descritti è avvenuta attraverso una metodologia progressiva di selezione, che ha progressivamente ristretto il campo delle opzioni:

\begin{itemize}
    \item \textbf{Studio individuale e discussione collettiva:} la fase iniziale ha previsto una lettura e un'analisi autonoma dei singoli capitolati da parte di ogni membro. A questo è seguita una prima valutazione di gruppo, fondamentale per confrontare le opinioni, discutere i punti di forza e di debolezza di ogni proposta e far emergere le preferenze condivise del team.
    \item \textbf{Colloqui con le aziende e decisione finale:} sulla base della scrematura iniziale, il gruppo ha organizzato degli incontri mirati con le aziende proponenti dei progetti verso cui si sentiva maggiormente orientato. Questi confronti diretti hanno permesso di verificare la disponibilità dei referenti e risolvere le ultime incertezze, guidando il team verso una decisione finale pienamente informata e concorde.
\end{itemize}

\section{Analisi del capitolato selezionato}
\subsection{C5 - Nexum}

Il capitolato C5, proposto da Eggon S.r.l., è stato selezionato dal gruppo per il suo forte equilibrio tra interesse tecnologico, valore formativo e aderenza a un contesto industriale concreto. La proposta si inserisce all'interno di \textit{NEXUM}, una piattaforma HR già esistente, e prevede l'estensione del prodotto con moduli basati su intelligenza artificiale generativa, automazione documentale e integrazione con i processi dei Consulenti del Lavoro.

\subsubsection{Obiettivo}

L'obiettivo del progetto è potenziare la piattaforma \textit{NEXUM} con funzionalità intelligenti e interoperabili, in grado di migliorare i processi HR e il dialogo tra aziende, dipendenti e studi dei Consulenti del Lavoro. In particolare, il capitolato individua due direttrici principali di sviluppo: un \textbf{AI Assistant Generativo} per la creazione di contenuti aziendali e un \textbf{AI Co-Pilot} per il trattamento, il riconoscimento e il dispaccio automatico dei documenti gestiti dai CdL.

In termini operativi, il progetto richiede la realizzazione di componenti che non siano semplici prototipi isolati, ma moduli effettivamente integrabili nell'architettura esistente di \textit{NEXUM}. Il valore della proposta risiede proprio nella possibilità di lavorare su funzionalità orientate a bisogni aziendali reali, con attenzione sia agli aspetti applicativi sia a quelli di qualità, sicurezza e tracciabilità.

\subsubsection{Caratteristiche funzionali principali}

Dal punto di vista funzionale, il capitolato propone uno scenario ampio e ben strutturato. Un primo ambito riguarda l'\textbf{AI Assistant Generativo}, pensato per supportare gli utenti della dashboard amministrativa nella redazione di contenuti aziendali. Le funzionalità centrali comprendono:

\begin{itemize}
    \item generazione automatica di titolo, contenuto e immagine di copertina a partire da un prompt testuale;
    \item personalizzazione del tono e dello stile comunicativo in base al contesto aziendale;
    \item salvataggio e riuso dei prompt precedenti, con storico consultabile;
    \item sistema di rating per valutare la qualità dei risultati prodotti dall'AI;
    \item revisione, modifica manuale e pubblicazione del contenuto direttamente nella piattaforma;
    \item raccolta di analytics e reportistica sull'utilizzo del modulo generativo.
\end{itemize}

Un secondo ambito riguarda invece l'\textbf{AI Co-Pilot per i Consulenti del Lavoro}, orientato all'automazione del flusso documentale. In questo caso il capitolato richiede:

\begin{itemize}
    \item caricamento di documenti nel repository CdL attraverso dashboard amministrativa;
    \item classificazione automatica dei documenti tramite riconoscimento e OCR;
    \item estrazione e riconoscimento dei destinatari mediante tecniche di matching e disambiguazione;
    \item suddivisione di documenti massivi per singolo destinatario;
    \item revisione \textit{human-in-the-loop} in caso di bassa confidenza o ambiguità;
    \item generazione automatica della lista di distribuzione e del messaggio di dispaccio;
    \item invio, tracciamento, prova di consegna, ricerca, audit e conservazione dei documenti;
    \item gestione di ruoli, permessi, policy di retention e monitoraggio continuo della qualità AI.
\end{itemize}

Nel complesso, il capitolato si distingue per l'ampiezza del dominio coperto: non si limita a una singola funzionalità innovativa, ma combina AI generativa, OCR, workflow documentali, sicurezza, analytics e integrazione con processi amministrativi concreti. Questa varietà è stata percepita come uno dei principali punti di interesse del progetto, poiché consente di affrontare problematiche diverse ma tra loro coerenti.

\subsubsection{Requisiti obbligatori e opzionali}

Dal punto di vista della pianificazione, il capitolato non presenta una distinzione rigida tra requisiti obbligatori e opzionali analoga a quella osservabile in altri casi. La struttura è infatti organizzata soprattutto attorno a requisiti di business, casi d'uso, vincoli, criteri di accettazione e deliverable. Per il gruppo, ciò significa che il nucleo obbligatorio risulta particolarmente ampio e comprende sia la realizzazione dei due moduli principali sia la loro integrazione con l'ecosistema \textit{NEXUM}.

Tra i requisiti da considerare essenziali rientrano:

\begin{itemize}
    \item integrazione con backend, frontend e database della piattaforma esistente;
    \item disponibilità delle funzionalità su dashboard amministrativa e, dove previsto, su PWA;
    \item implementazione dei principali casi d'uso dell'AI Assistant Generativo;
    \item implementazione del flusso documentale dell'AI Co-Pilot, dal caricamento al dispatch;
    \item presenza di meccanismi di audit trail, tracciabilità, ruoli e permessi;
    \item produzione di deliverable concreti, demo funzionante e documentazione tecnica.
\end{itemize}

Più che veri requisiti opzionali, il capitolato introduce elementi di qualità e di maturità del prodotto che ne aumentano il valore complessivo, come metriche di performance, monitoraggio dei KPI, miglioramento continuo dei modelli AI, attenzione alla compliance e possibile evoluzione progressiva dei moduli. Questa impostazione è stata valutata positivamente perché favorisce una crescita incrementale del sistema, pur mantenendo obiettivi progettuali chiari e industrialmente rilevanti.

\subsubsection{Considerazioni tecnologiche}

Uno dei punti di forza del capitolato è la chiarezza dello stack tecnologico e architetturale. La documentazione indica infatti un ecosistema moderno e realistico, composto da \textbf{Angular} per la dashboard amministrativa, \textbf{Next.js} per la PWA, \textbf{Ruby on Rails} in modalità API-first per il backend, \textbf{PostgreSQL} per la persistenza, \textbf{Redis} e \textbf{Sidekiq} per code e processi asincroni, e un'infrastruttura \textbf{AWS} articolata su servizi come ECS Fargate, S3, CloudFront, Cognito, SQS, SES, CloudWatch, GuardDuty e WAF.

Questa chiarezza tecnologica è stata valutata molto positivamente perché riduce l'ambiguità iniziale e rende più concreto il lavoro di analisi architetturale. Allo stesso tempo, la varietà delle tecnologie coinvolte aumenta la complessità del progetto: il gruppo deve confrontarsi con sviluppo web full stack, integrazione cloud, sicurezza, gestione documentale, servizi asincroni e componenti AI, il tutto in un contesto di prodotto reale e non puramente sperimentale.

\subsubsection{Incontro con il proponente}

Dall'incontro con il proponente sono emersi i seguenti chiarimenti:
\begin{itemize}
    \item \textbf{Collaborazione:} incontri prevalentemente telematici, con cadenza periodica definita dal team;
    \item \textbf{Perimetro del progetto:} sviluppo di una prova tecnica autonoma, senza integrazione diretta nei sistemi aziendali esistenti;
    \item \textbf{Indicazioni tecnologiche:} consigliato ecosistema AWS (es. Bedrock/Comprehend), con piena libertà di stack e preferenza per tecnologie già note al gruppo;
    \item \textbf{Approccio progettuale:} uso dell'AI motivato dall'eterogeneità documentale, con impostazione formativa e possibile ricalibrazione di tempi e attività.
\end{itemize}

Per ulteriori informazioni, consultare il \href{https://alittlebyte-19.github.io/Documentazione/candidatura/verbali/Verbali%20Esterni/Verbale_Esterno_2026-03-16.pdf}{\textcolor{blue}{\underline{verbale dell'incontro}}}.

\subsubsection{Valutazione del gruppo}

Durante la valutazione, il gruppo ha riconosciuto nel capitolato C5 diversi punti di forza:

\begin{itemize}
    \item forte aderenza a un caso industriale reale, con possibilità di produrre risultati concretamente integrabili nel prodotto aziendale;
    \item equilibrio tra innovazione e applicabilità pratica, grazie alla combinazione di AI generativa e automazione documentale;
    \item chiarezza dei casi d'uso, dei vincoli, dei criteri di accettazione e dei deliverable attesi;
    \item stack moderno e formativo, in grado di offrire esperienza su tecnologie attuali di mercato;
    \item opportunità di lavorare secondo metodologie Agile e con strumenti professionali in un contesto collaborativo strutturato.
\end{itemize}

Parallelamente, il gruppo ha individuato alcune criticità:

\begin{itemize}
    \item ampiezza funzionale elevata, che richiede una pianificazione attenta per evitare dispersione dello sforzo;
    \item complessità architetturale non trascurabile, dovuta al numero di servizi e integrazioni coinvolte;
    \item necessità di affrontare temi delicati come sicurezza, GDPR, auditabilità e affidabilità dei sistemi AI.
\end{itemize}

Nonostante tali criticità, il gruppo ha individuato con convinzione nel capitolato C5 la scelta migliore. La decisione è maturata in modo solido e condiviso, poiché C5 è risultato il più allineato agli obiettivi formativi e progettuali del team: combina interesse tecnologico, struttura chiara, forte rilevanza industriale e una concreta sostenibilità realizzativa entro i vincoli temporali del corso.

\section{Analisi degli altri capitolati}
\subsection{C1 - Automated EN18031 Compliance Verification}

Il capitolato C1, proposto da Bluewind S.r.l., è stato inizialmente considerato dal gruppo per la sua aderenza a un tema attuale e rilevante: la conformità alla norma EN 18031, armonizzata rispetto alla direttiva RED per i dispositivi radio immessi nel mercato europeo.


\subsubsection{Obiettivo}

L'obiettivo del progetto è la realizzazione di un'applicazione in grado di supportare e automatizzare la verifica di conformità ai requisiti EN 18031 tramite esecuzione guidata di Decision Tree. Per ogni requisito analizzato, il sistema deve produrre un esito tra \textit{Pass}, \textit{Fail} e \textit{Not Applicable}, riducendo il carico di lavoro manuale e migliorando la tracciabilità delle decisioni.

In termini operativi, il software deve guidare l'utente nell'attraversamento dei nodi decisionali (domande con risposta \textit{Yes/No}) fino al nodo terminale e registrare in modo coerente percorso seguito, motivazioni e risultato finale per ciascun requisito.

\subsubsection{Caratteristiche funzionali principali}

Il capitolato richiede in particolare:

\begin{itemize}
    \item importazione di documenti tecnici e file esterni descrittivi dei Decision Tree (ad esempio XML e JSON);
    \item esecuzione interattiva degli alberi decisionali nel rispetto delle dipendenze gerarchiche tra requisiti;
    \item visualizzazione aggregata dei risultati tramite dashboard;
    \item possibilità di modifica e aggiornamento dei Decision Tree, con salvataggio in formati standard.
\end{itemize}

Oltre a queste funzionalità, il capitolato evidenzia anche l'esigenza di standardizzare la produzione della documentazione tecnica derivante dalla valutazione, così da rendere più semplice il riesame delle decisioni e il loro aggiornamento in caso di revisione normativa.

Il caso studio proposto (macchina del caffè connessa via Wi-Fi con comunicazione MQTT) rende il contesto concreto e utile per validare il comportamento dell'applicazione sui requisiti ACM e AUM. La presenza di un dominio applicativo reale costituisce un punto di forza, poiché permette di testare il flusso end-to-end su uno scenario non artificiale.

\subsubsection{Requisiti obbligatori e opzionali}

Dal punto di vista della pianificazione, il capitolato presenta una distinzione chiara tra requisiti obbligatori e opzionali.

Tra i requisiti obbligatori risultano centrali l'importazione di formati standard, l'esecuzione automatizzata dei Decision Tree, il rispetto delle dipendenze gerarchiche e la presenza di una dashboard per la visualizzazione dei risultati. Si tratta di una base funzionale ben definita, orientata a garantire tracciabilità, correttezza del flusso decisionale e leggibilità degli esiti.

Tra i requisiti opzionali emergono invece funzionalità di maggiore valore aggiunto, che possono incrementare qualità del prodotto ma anche complessità implementativa:

\begin{itemize}
    \item editor integrato per la modifica dei Decision Tree;
    \item esportazione dei risultati e dei dati in più formati;
    \item giustificazione esplicita dei risultati ottenuti;
    \item editor grafico avanzato per la modellazione degli alberi decisionali.
\end{itemize}

Questa separazione è stata valutata positivamente dal gruppo perché facilita la definizione di una baseline minima e consente un'estensione incrementale del progetto.

\subsubsection{Considerazioni tecnologiche}

Il proponente non impone vincoli rigidi su stack, architettura o tipologia applicativa (desktop o web), lasciando libertà progettuale al team. L'unica preferenza esplicita riguarda l'uso di Python 3.x lato backend, preferibilmente con gestione del progetto tramite strumenti di packaging moderni.

Questa flessibilità rappresenta contemporaneamente un vantaggio e un costo: da un lato permette al gruppo di scegliere soluzioni aderenti alle proprie competenze, dall'altro richiede una fase preliminare più ampia di analisi e confronto tecnico prima di convergere su un'architettura definitiva.

\subsubsection{Incontro con il proponente}

Dall'incontro con il proponente sono emersi i seguenti chiarimenti:
\begin{itemize}
    \item \textbf{Tecnologie:} piena libertà nello stack frontend; Python suggerito lato backend per facilitare il supporto aziendale;
    \item \textbf{Persistenza dati:} database non obbligatorio, è accettato anche il salvataggio locale su file;
    \item \textbf{Architettura:} consigliato caricamento dinamico degli alberi decisionali (JSON/XML) e gestione automatica delle dipendenze gerarchiche;
    \item \textbf{Collaborazione:} incontri a richiesta, con contatti rapidi via email/Telegram e call quando necessario.
\end{itemize}

Per ulteriori informazioni, consultare il \href{https://alittlebyte-19.github.io/Documentazione/candidatura/verbali/Verbali%20Esterni/Verbale_Esterno_2026-03-17.pdf}{\textcolor{blue}{\underline{verbale dell'incontro}}}.

\subsubsection{Valutazione del gruppo}

Durante la valutazione, il gruppo ha riconosciuto i principali punti di forza del capitolato:

\begin{itemize}
    \item requisiti chiari e ben strutturati;
    \item caso studio concreto e rilevanza industriale;
    \item buona disponibilità del proponente al supporto.
\end{itemize}

Sono emerse però anche alcune criticità:

\begin{itemize}
    \item curva di apprendimento iniziale elevata legata all'interpretazione della norma;
    \item minore margine creativo rispetto ad altri capitolati.
\end{itemize}

In conclusione, C1 è stato ritenuto solido e ben definito, ma non selezionato perché meno allineato agli interessi e agli obiettivi formativi del gruppo rispetto al capitolato scelto.

\subsection{C2 - Code Guardian}

Il capitolato C2, proposto da Var Group S.p.A., presenta lo sviluppo di una piattaforma web orientata all'audit dei repository software, con particolare attenzione a qualità del codice, sicurezza e manutenibilità.

\subsubsection{Obiettivo}

L'obiettivo del progetto è realizzare un sistema modulare, basato su un'architettura ad agenti coordinati da un orchestratore centrale, in grado di analizzare repository GitHub e produrre report automatici sullo stato complessivo del progetto.

La piattaforma deve non solo individuare criticità tecniche e di sicurezza, ma anche suggerire azioni correttive (\textit{remediation}) per migliorare progressivamente qualità, robustezza e mantenibilità del software analizzato.

\subsubsection{Caratteristiche funzionali principali}

Dal punto di vista funzionale, il capitolato richiede in particolare:

\begin{itemize}
    \item analisi automatica dei repository per identificare stack tecnologico, dipendenze e versioni utilizzate;
    \item verifica della presenza e della qualità dei test, con attenzione alla copertura e all'affidabilità complessiva della suite;
    \item audit di sicurezza orientato alle linee guida OWASP, con evidenziazione di vulnerabilità e configurazioni rischiose;
    \item valutazione della qualità della documentazione tecnica (README, API docs, guide operative);
    \item dashboard web per visualizzare stato generale, criticità e priorità di intervento.
\end{itemize}

Un elemento distintivo del capitolato è la componente di remediation: il sistema non si limita a segnalare problemi, ma propone interventi migliorativi concreti, favorendo un approccio operativo e continuo alla qualità del codice.

\subsubsection{Requisiti obbligatori e opzionali}

La documentazione propone un percorso strutturato che include attività iniziali di analisi e progettazione (raccolta requisiti, user story mapping, modellazione UML e schema dati), seguite dalla realizzazione di un MVP funzionante con demo e documentazione tecnica.

Dal punto di vista del gruppo, questa impostazione è stata considerata positiva perché orientata a un processo realistico di sviluppo, ma al tempo stesso impegnativa in termini di organizzazione e carico progettuale. I requisiti obbligatori non si limitano infatti alle funzionalità applicative finali, ma comprendono anche una parte iniziale significativa di analisi, progettazione e documentazione.

Tra le estensioni possibili, risultano particolarmente interessanti:

\begin{itemize}
    \item integrazione CI/CD più avanzata;
    \item analisi storica delle versioni del repository;
    \item sistemi di ranking e confronto tra progetti;
    \item meccanismi interattivi di notifica e remediation.
\end{itemize}

\subsubsection{Considerazioni tecnologiche}

Il capitolato indica con chiarezza lo stack di riferimento: Node.js e Python per backend e orchestrazione, React.js per il frontend, MongoDB/PostgreSQL per la persistenza, GitHub Actions per l'automazione e AWS per l'infrastruttura cloud.

La chiarezza tecnologica è stata valutata favorevolmente perché riduce l'ambiguità iniziale e consente una pianificazione più precisa. Tuttavia, l'eterogeneità dello stack richiede competenze distribuite su più aree (sviluppo web, automazione, cloud, sicurezza), aumentando la complessità complessiva del progetto.

\subsubsection{Valutazione del gruppo}

Il gruppo ha riconosciuto in C2 diversi punti di forza:

\begin{itemize}
    \item tema attuale e di forte utilità pratica nel contesto industriale;
    \item architettura moderna e potenzialmente molto formativa;
    \item obiettivi ben definiti e tecnologie chiaramente esplicitate.
\end{itemize}

Parallelamente, sono emerse alcune criticità rilevanti:

\begin{itemize}
    \item complessità tecnica elevata dovuta alla combinazione di multi-agente, sicurezza applicativa e integrazioni esterne;
    \item carico consistente legato ai requisiti di qualità e validazione del software;
    \item interesse complessivo del gruppo inferiore rispetto ad altri capitolati ritenuti più allineati alle preferenze progettuali condivise.
\end{itemize}

In conclusione, C2 è stato considerato un capitolato solido, moderno e professionalmente molto valido, ma non è stato scelto dal gruppo aLittleByte poiché valutato meno aderente all'equilibrio desiderato tra interesse, complessità e obiettivi formativi.

\subsection{C4 - L'app che Protegge e Trasforma}

Il capitolato C4, proposto da Miriade S.r.l., è stato preso in considerazione dal gruppo per il suo tema attualmente molto discusso e impattante: la violenza di genere e la necessità di mettere a disposizione delle persone un mezzo tangibile, affidabile e all'avanguardia in termini tecnologici per le attività di prevenzione e di sostegno alle vittime.

\subsubsection{Obiettivo}

L'obiettivo principale è fornire un'applicazione per sistemi operativi iOS e Android attiva, intelligente e sicura, in grado di riconoscere segnali di pericolo e offrire risorse concrete per l'autonomia e la sicurezza degli utenti. L’applicazione mobile, inoltre, deve essere un prodotto eticamente responsabile e ad impatto sociale positivo. In termini operativi, l’app deve identificare situazioni di rischio, inviare alert in modo discreto e garantire la totale privacy e sicurezza dell'utente, fornendo al contempo percorsi educativi e reti di supporto.

\subsubsection{Caratteristiche funzionali principali}

Le funzionalità principali richieste dall’azienda sono:
\begin{itemize}
    \item \textbf{Rilevamento e Alert:} il software aiuta l’utente ad individuare ``red flag'', ovvero comportamenti a rischio, tramite intelligenza artificiale con possibilità di inviare notifiche di allerta silenziose a reti di contatti fiduciari.
    \item \textbf{Risorse e supporto:} l’app supporta l’utente legalmente e psicologicamente su come uscire da relazioni violente e accedere a una rete di servizi geolocalizzati, integrando servizi di invio rapido di richieste di aiuto a numeri di emergenza e centri antiviolenza.
    \item \textbf{Funzionalità di Sicurezza Personalizzate:} modalità di personalizzazione dell’utente per nascondere l’app, creare un diario criptato e pianificare percorsi sicuri o di fuga con mappe interattive.
    \item \textbf{Formazione e Prevenzione:} moduli educativi interattivi per sensibilizzare sul tema della violenza di genere.
    \item \textbf{Community di Supporto:} uno spazio sicuro, con moderazione professionale, per la condivisione di esperienze e il supporto reciproco tra utenti.
\end{itemize}

Nel complesso, il capitolato delinea un prodotto mobile che non si limita alla gestione dell'emergenza, ma integra prevenzione, accompagnamento e autonomia personale. Il valore della proposta risiede proprio nell'unione tra componenti di riconoscimento del rischio, accesso immediato ai servizi territoriali e strumenti di protezione discreti, pensati per un contesto d'uso delicato e potenzialmente critico.

\subsubsection{Requisiti obbligatori e opzionali}

Dal punto di vista della pianificazione, i requisiti obbligatori comprendono il riconoscimento di segnali di rischio, l'invio di alert silenziosi, l'accesso a risorse di supporto geolocalizzate, la disponibilità di funzionalità di sicurezza personalizzate, la presenza di contenuti formativi e la creazione di uno spazio protetto di community. Si tratta quindi di una baseline già molto ampia, che combina aspetti applicativi, sociali e di protezione dei dati.

Accanto a questo nucleo minimo, il capitolato propone diverse estensioni opzionali che aumentano la completezza del prodotto e ne rafforzano l'impatto pratico:
\begin{itemize}
    \item integrazione con dispositivi indossabili, come smartwatch, per inviare allarmi o registrare eventi in modo ancora più discreto e immediato;
    \item modalità di ``check-in'' periodico, in cui il mancato riscontro dell'utente attiva automaticamente una segnalazione verso i contatti di emergenza;
    \item funzionalità di ``fuga rapida'', con attivazione combinata di allarme, registrazione audio-video, invio posizione e messaggio predefinito;
    \item rilevamento ambientale tramite analisi di suoni e rumori associati a situazioni di pericolo, previo consenso dell'utente;
    \item gamification dei moduli educativi, con punti, livelli e ricompense per aumentare coinvolgimento ed efficacia preventiva;
    \item supporto multilingue avanzato, esteso anche a lingue meno diffuse e dialetti;
    \item reportistica anonima per enti di ricerca, con dati anonimizzati condivisibili solo previo consenso esplicito;
    \item integrazione con la legislazione locale e internazionale, per fornire informazioni aggiornate e guide pratiche su come agire;
    \item accessibilità avanzata per persone con disabilità visive, uditive o motorie.
\end{itemize}

Questa distinzione è stata letta dal gruppo come segnale di un capitolato ricco ma impegnativo: le funzionalità opzionali non sono semplici rifiniture, bensì estensioni di grande valore che aumentano notevolmente ampiezza, responsabilità progettuale e complessità realizzativa.

\subsubsection{Considerazioni tecnologiche}

Il proponente fornisce indicazioni architetturali chiare. Per lo sviluppo mobile, si raccomanda vivamente l'adozione del framework \textbf{Flutter} al fine di ottimizzare i tempi di sviluppo e la manutenibilità sia su piattaforme iOS che Android. Per il Backend, si propone l'adozione di un'architettura \textbf{Serverless}, con hosting preferibilmente su \textbf{AWS}. La soluzione suggerita contempla l'impiego di Amazon API Gateway, AWS Lambda per l'implementazione dei microservizi, e l'utilizzo di database sia NoSQL (Amazon DynamoDB) che relazionali (Amazon RDS). È inoltre richiesta l'integrazione con Amazon SageMaker e Bedrock al fine di gestire i modelli di Machine Learning e l'analisi del linguaggio naturale.

\subsubsection{Valutazione del gruppo}

Durante la valutazione, il gruppo ha riconosciuto i punti di forza del capitolato: l’importanza sociale e la grande disponibilità del proponente. Sono stati inoltre considerati positivamente:
\begin{itemize}
    \item il supporto completo da parte dell’azienda: Miriade S.r.l. offre referenti di supporto durante tutto il ciclo di vita dello sviluppo software;
    \item la possibilità di utilizzare un ambiente della suite ``Atlassian'' (Jira, Bitbucket) per facilitare lo sviluppo;
    \item la formazione iniziale sulla prevenzione della violenza di genere e sul supporto alle vittime.
\end{itemize}

Nonostante ciò, il progetto non è stato selezionato. Le principali motivazioni sono state: la necessità di un investimento di tempo iniziale troppo gravoso per il gruppo, la preferenza del team verso domini meno critici dal punto di vista etico e le implementazioni troppo impegnative riguardanti la sicurezza dei dati. In conclusione, la valutazione ha rilevato una minore aderenza alle tempistiche accademiche e alle conoscenze tecnologiche del gruppo rispetto al capitolato scelto.

\subsection{C6 - Second Brain}

Il capitolato C6, proposto da Zucchetti S.p.A., è risultato un punto di interesse per il gruppo per il suo forte orientamento verso le tecnologie emergenti, in particolare l’integrazione di modelli \textbf{LLM (Large Language Models)} nel lavoro quotidiano.

\subsubsection{Obiettivo}

L’obiettivo consiste nel creare una Web App di un editor di testo basato sul linguaggio di Markup \textbf{Markdown} con integrazione di intelligenza artificiale. L’applicazione dovrà eseguire operazioni di riassunto, critica di testo, traduzione e riformulazione interagendo con gli LLM e dovrà supportare prompt che permettano di generare testi in modo automatico.

\subsubsection{Caratteristiche funzionali principali}

Il nucleo obbligatorio del capitolato ruota attorno a un editor Markdown arricchito dalle capacità di un LLM, con l'obiettivo di supportare l'utente nella scrittura, nella rielaborazione e nella generazione di contenuti testuali. In particolare, la piattaforma deve consentire:

\begin{itemize}
    \item Editing di testo in un'area che accetta testo e marcatori in formato Markdown;
    \item visualizzazione grafica del testo applicando lo stile corrispondente ai marcatori;
    \item accesso a un LLM che possa operare sul testo per intero o per parti;
    \item implementazione di comandi base di riassunto, riscrittura e traduzione;
    \item implementazione di comandi di critica basati sulla ``tecnica dei sei cappelli'' di Edward De Bono;
    \item gestione di un prompt associato per generare testi in maniera automatica;
    \item persistenza e lettura dei dati delle note in formato testuale.
\end{itemize}

Il capitolato chiarisce anche il significato operativo di queste funzioni: l'editor deve essere semplice e immediato, la preview deve rendere leggibile il risultato finale della scrittura in Markdown, mentre l'interazione con l'LLM deve poter avvenire sia sull'intero testo sia su sue porzioni selezionate. Particolare rilievo assume la progettazione dei prompt, soprattutto per le operazioni di critica e riscrittura, considerate il vero elemento di ricerca del progetto.

\subsubsection{Requisiti obbligatori e opzionali}

I requisiti obbligatori comprendono quindi non solo la realizzazione dell'editor e della preview, ma anche l'integrazione di operazioni testuali fondamentali --- riassunto, riscrittura, traduzione, critica strutturata e generazione guidata da prompt --- oltre al salvataggio e recupero delle note in formato testuale. Il capitolato sottolinea inoltre che questi requisiti minimi devono essere sviluppati secondo una progressione chiara, passando da una fase iniziale di \textit{Proof of Concept} a un \textit{Minimum Viable Product} completo, con crescente attenzione a usabilità, robustezza e test automatici.

Sul piano delle estensioni, Zucchetti S.p.A. propone di sfruttare ulteriormente la potenza degli LLM tramite:
\begin{itemize}
    \item archivio delle note in un database server-side;
    \item possibilità di recuperare e modificare una nota scritta in precedenza;
    \item associazione delle note tramite link integrati nel markup;
    \item utilizzo delle note collegate durante la creazione o modifica dei testi nell’editor.
\end{itemize}

Questi requisiti opzionali introducono una vera componente di knowledge management: le note non restano documenti isolati, ma diventano una base informativa collegata e riusabile, capace di alimentare l'LLM con contesto aggiuntivo. È proprio questa estensione a dare pieno significato al nome ``Second Brain'', trasformando l'editor da semplice strumento di scrittura assistita a sistema di supporto alla riflessione e al recupero strutturato delle idee.

\subsubsection{Considerazioni tecnologiche}

Il proponente lascia ampia libertà nella scelta dello stack tecnologico, suggerendo l'uso di HTML5 per il frontend e librerie Open Source per il rendering. Per l'interazione con l'IA, viene fornito l'uso di API compatibili con lo standard OpenAI (Gemini, Mistral o Gemma). Qualora si desideri sviluppare il database, l'azienda propone la gestione delle chiamate HTTP tramite un backend che può essere realizzato in \textbf{Java} o \textbf{Python}.

\subsubsection{Incontro con il proponente}

Dall'incontro con il proponente sono emersi i seguenti chiarimenti:
\begin{itemize}
    \item \textbf{Frequenza incontri:} a discrezione del gruppo, previa organizzazione;
    \item \textbf{Tecnologie:} scelta libera sui framework; per il backend è consigliato Java ma non è vincolante;
    \item \textbf{Modelli LLM:} è possibile utilizzare modelli OpenAI o API key fornite per Gemini, Mistral o Gemma;
    \item \textbf{Storage:} è consigliato l'uso di \textbf{PostgreSQL}, mentre è sconsigliato l'utilizzo di database NoSQL.
\end{itemize}

Per ulteriori informazioni, consultare il \href{https://alittlebyte-19.github.io/Documentazione/candidatura/verbali/Verbali%20Esterni/Verbale_Esterno_2026-03-11.pdf}{\textcolor{blue}{\underline{verbale dell'incontro}}}.

\subsubsection{Valutazione del gruppo}

Il gruppo ha riconosciuto i seguenti punti di forza:
\begin{itemize}
    \item architettura flessibile con possibilità di progettazione libera;
    \item disponibilità del proponente a garantire l'accesso diretto ai modelli LLM;
    \item tempistiche di realizzazione ottimali grazie al focus sul prompt engineering.
\end{itemize}

I punti di debolezza risiedono nelle scarse specifiche vincolanti, che hanno ridotto l'interesse nello studio di nuovi framework frontend. In conclusione, il capitolato C6 è risultato un valido concorrente poiché ne permette la realizzazione in tempi ragionevoli, ma l’interesse generale non è stato sufficiente per la maggioranza dei membri rispetto al capitolato scelto.

\end{document}
